% ============================================================================
% Abstract
% % ============================================================================
% Abstract
% % ============================================================================
% Abstract
% % ============================================================================
% Abstract
% \input{abstract} inside \begin{abstract} ... \end{abstract}
% ============================================================================

Masked language models pretrained on SMILES offer a route to de novo molecular
design that needs no task-specific generative training: masking selected tokens
of a known ligand and resampling them reduces generation to local, token-level
editing of an existing scaffold. We ask whether the choice of \emph{which}
tokens to mask carries structural information, by masking the atoms that PLIP
identifies as participating in protein--ligand contacts in deposited
co-crystal structures and comparing against uniformly random masking of the
same parent molecules at the same rate, so that the two arms differ only in
mask placement. Because the pretrained model infills such masks with low
chemical fidelity, we fine-tune ChemBERTa against a composite, non-differentiable
objective over RDKit validity, drug-likeness, synthetic accessibility,
divergence from the parent scaffold, PAINS/Brenk structural alerts, and
predicted Tox21 activity from a multitask classifier fine-tuned on the same
backbone. Two estimators are implemented as a controlled comparison --- a
REINFORCE policy gradient with LoRA adapters and an exact KL anchor to the
pretrained policy, and best-of-$K$ selection followed by ordinary
cross-entropy --- with the training pairs, masking rate, rollout distribution
and scoring code held fixed between them.

Initial results are encouraging but preliminary. Across three epochs of
policy-gradient fine-tuning the fraction of infilled molecules that parse rose
from 20.9\% to 22.7\% and the mean composite reward from 0.101 to 0.110, with
the drug-likeness, accessibility and novelty terms improving in step and the
policy held near 2 nats per masked position from its pretrained distribution.
Absolute validity remains low, and the gain is modest against the range the
objective defines.

Part of that ceiling appears attributable to the composition of the training
corpus rather than to the objective. PDB-derived ligand sets are dominated by
a few chemotypes: the acyclic, Fe-porphyrin and Mg-porphyrin frameworks alone
account for 28\% of distinct parent molecules, and crystallisation additives
--- sulfate and glycerol --- for a quarter of all deposited instances. We
therefore deduplicate to one pair per distinct parent and evaluate on a
Bemis--Murcko framework partition, under which no held-out molecule shares a
ring system with any training molecule, reporting training and held-out
distributions separately. Ongoing work is directed at making the model
generalise across chemically diverse molecular groups rather than at improving
aggregate numbers on this skewed distribution.
 inside \begin{abstract} ... \end{abstract}
% ============================================================================

Masked language models pretrained on SMILES offer a route to de novo molecular
design that needs no task-specific generative training: masking selected tokens
of a known ligand and resampling them reduces generation to local, token-level
editing of an existing scaffold. We ask whether the choice of \emph{which}
tokens to mask carries structural information, by masking the atoms that PLIP
identifies as participating in protein--ligand contacts in deposited
co-crystal structures and comparing against uniformly random masking of the
same parent molecules at the same rate, so that the two arms differ only in
mask placement. Because the pretrained model infills such masks with low
chemical fidelity, we fine-tune ChemBERTa against a composite, non-differentiable
objective over RDKit validity, drug-likeness, synthetic accessibility,
divergence from the parent scaffold, PAINS/Brenk structural alerts, and
predicted Tox21 activity from a multitask classifier fine-tuned on the same
backbone. Two estimators are implemented as a controlled comparison --- a
REINFORCE policy gradient with LoRA adapters and an exact KL anchor to the
pretrained policy, and best-of-$K$ selection followed by ordinary
cross-entropy --- with the training pairs, masking rate, rollout distribution
and scoring code held fixed between them.

Initial results are encouraging but preliminary. Across three epochs of
policy-gradient fine-tuning the fraction of infilled molecules that parse rose
from 20.9\% to 22.7\% and the mean composite reward from 0.101 to 0.110, with
the drug-likeness, accessibility and novelty terms improving in step and the
policy held near 2 nats per masked position from its pretrained distribution.
Absolute validity remains low, and the gain is modest against the range the
objective defines.

Part of that ceiling appears attributable to the composition of the training
corpus rather than to the objective. PDB-derived ligand sets are dominated by
a few chemotypes: the acyclic, Fe-porphyrin and Mg-porphyrin frameworks alone
account for 28\% of distinct parent molecules, and crystallisation additives
--- sulfate and glycerol --- for a quarter of all deposited instances. We
therefore deduplicate to one pair per distinct parent and evaluate on a
Bemis--Murcko framework partition, under which no held-out molecule shares a
ring system with any training molecule, reporting training and held-out
distributions separately. Ongoing work is directed at making the model
generalise across chemically diverse molecular groups rather than at improving
aggregate numbers on this skewed distribution.
 inside \begin{abstract} ... \end{abstract}
% ============================================================================

Masked language models pretrained on SMILES offer a route to de novo molecular
design that needs no task-specific generative training: masking selected tokens
of a known ligand and resampling them reduces generation to local, token-level
editing of an existing scaffold. We ask whether the choice of \emph{which}
tokens to mask carries structural information, by masking the atoms that PLIP
identifies as participating in protein--ligand contacts in deposited
co-crystal structures and comparing against uniformly random masking of the
same parent molecules at the same rate, so that the two arms differ only in
mask placement. Because the pretrained model infills such masks with low
chemical fidelity, we fine-tune ChemBERTa against a composite, non-differentiable
objective over RDKit validity, drug-likeness, synthetic accessibility,
divergence from the parent scaffold, PAINS/Brenk structural alerts, and
predicted Tox21 activity from a multitask classifier fine-tuned on the same
backbone. Two estimators are implemented as a controlled comparison --- a
REINFORCE policy gradient with LoRA adapters and an exact KL anchor to the
pretrained policy, and best-of-$K$ selection followed by ordinary
cross-entropy --- with the training pairs, masking rate, rollout distribution
and scoring code held fixed between them.

Initial results are encouraging but preliminary. Across three epochs of
policy-gradient fine-tuning the fraction of infilled molecules that parse rose
from 20.9\% to 22.7\% and the mean composite reward from 0.101 to 0.110, with
the drug-likeness, accessibility and novelty terms improving in step and the
policy held near 2 nats per masked position from its pretrained distribution.
Absolute validity remains low, and the gain is modest against the range the
objective defines.

Part of that ceiling appears attributable to the composition of the training
corpus rather than to the objective. PDB-derived ligand sets are dominated by
a few chemotypes: the acyclic, Fe-porphyrin and Mg-porphyrin frameworks alone
account for 28\% of distinct parent molecules, and crystallisation additives
--- sulfate and glycerol --- for a quarter of all deposited instances. We
therefore deduplicate to one pair per distinct parent and evaluate on a
Bemis--Murcko framework partition, under which no held-out molecule shares a
ring system with any training molecule, reporting training and held-out
distributions separately. Ongoing work is directed at making the model
generalise across chemically diverse molecular groups rather than at improving
aggregate numbers on this skewed distribution.
 inside \begin{abstract} ... \end{abstract}
% ============================================================================

Masked language models pretrained on SMILES offer a route to de novo molecular
design that needs no task-specific generative training: masking selected tokens
of a known ligand and resampling them reduces generation to local, token-level
editing of an existing scaffold. We ask whether the choice of \emph{which}
tokens to mask carries structural information, by masking the atoms that PLIP
identifies as participating in protein--ligand contacts in deposited
co-crystal structures and comparing against uniformly random masking of the
same parent molecules at the same rate, so that the two arms differ only in
mask placement. Because the pretrained model infills such masks with low
chemical fidelity, we fine-tune ChemBERTa against a composite, non-differentiable
objective over RDKit validity, drug-likeness, synthetic accessibility,
divergence from the parent scaffold, PAINS/Brenk structural alerts, and
predicted Tox21 activity from a multitask classifier fine-tuned on the same
backbone. Two estimators are implemented as a controlled comparison --- a
REINFORCE policy gradient with LoRA adapters and an exact KL anchor to the
pretrained policy, and best-of-$K$ selection followed by ordinary
cross-entropy --- with the training pairs, masking rate, rollout distribution
and scoring code held fixed between them.

Initial results are encouraging but preliminary. Across three epochs of
policy-gradient fine-tuning the fraction of infilled molecules that parse rose
from 20.9\% to 22.7\% and the mean composite reward from 0.101 to 0.110, with
the drug-likeness, accessibility and novelty terms improving in step and the
policy held near 2 nats per masked position from its pretrained distribution.
Absolute validity remains low, and the gain is modest against the range the
objective defines.

Part of that ceiling appears attributable to the composition of the training
corpus rather than to the objective. PDB-derived ligand sets are dominated by
a few chemotypes: the acyclic, Fe-porphyrin and Mg-porphyrin frameworks alone
account for 28\% of distinct parent molecules, and crystallisation additives
--- sulfate and glycerol --- for a quarter of all deposited instances. We
therefore deduplicate to one pair per distinct parent and evaluate on a
Bemis--Murcko framework partition, under which no held-out molecule shares a
ring system with any training molecule, reporting training and held-out
distributions separately. Ongoing work is directed at making the model
generalise across chemically diverse molecular groups rather than at improving
aggregate numbers on this skewed distribution.
